% Disproving Bandeira's Matrix AM-GM Open Problem 0.2(a):
% An Exact Rational Certificate and Formal Verification
%
% Compile: latexmk -pdf -interaction=nonstopmode main.tex

\begin{filecontents*}{refs.bib}
@inproceedings{RechtRe2012,
  author    = {Benjamin Recht and Christopher R{\'e}},
  title     = {Beneath the Valley of the Noncommutative Arithmetic--Geometric Mean Inequality: Conjectures, Case-studies, and Consequences},
  booktitle = {Proceedings of the 25th Annual Conference on Learning Theory},
  series    = {PMLR},
  volume    = {23},
  pages     = {11.1--11.24},
  year      = {2012}
}
@inproceedings{LaiLim2020,
  author    = {Zehua Lai and Lek-Heng Lim},
  title     = {Recht--R{\'e} Noncommutative Arithmetic--Geometric Mean Conjecture is False},
  booktitle = {Proceedings of the 37th International Conference on Machine Learning},
  series    = {PMLR},
  volume    = {119},
  pages     = {5608--5617},
  year      = {2020}
}
@misc{Bandeira2015,
  author       = {Afonso S. Bandeira},
  title        = {Ten Lectures and Forty-Two Open Problems in the Mathematics of Data Science},
  howpublished = {MIT 18.S096 Lecture Notes},
  year         = {2015},
  note         = {Open Problem 0.2}
}
@article{Helton2002,
  author  = {J. William Helton},
  title   = {{``Positive''} noncommutative polynomials are sums of squares},
  journal = {Annals of Mathematics},
  volume  = {156},
  number  = {2},
  pages   = {675--694},
  year    = {2002}
}
@article{IsraelKrahmerWard2016,
  author  = {Arie Israel and Felix Krahmer and Rachel Ward},
  title   = {An arithmetic--geometric mean inequality for products of three matrices},
  journal = {Linear Algebra and its Applications},
  volume  = {488},
  pages   = {1--12},
  year    = {2016}
}
@article{Sturm1999,
  author  = {Jos F. Sturm},
  title   = {Using {SeDuMi} 1.02, a {MATLAB} toolbox for optimization over symmetric cones},
  journal = {Optimization Methods and Software},
  volume  = {11--12},
  pages   = {625--653},
  year    = {1999}
}
@inproceedings{moura2021lean4,
  author    = {Leonardo de Moura and Sebastian Ullrich},
  title     = {The {Lean}~4 Theorem Prover and Programming Language},
  booktitle = {Automated Deduction -- {CADE} 28},
  pages     = {625--635},
  year      = {2021}
}
@misc{mathlib4,
  author       = {{Mathlib Community}},
  title        = {Mathlib4: the math library for {Lean}~4},
  howpublished = {\url{https://github.com/leanprover-community/mathlib4}},
  year         = {2024}
}
@article{Lasserre2001,
  author  = {Jean B. Lasserre},
  title   = {Global optimization with polynomials and the problem of moments},
  journal = {SIAM Journal on Optimization},
  volume  = {11},
  number  = {3},
  pages   = {796--817},
  year    = {2001}
}
@article{Fujii2016,
  author  = {Jun Ichi Fujii and Masatoshi Fujii and Yuki Seo and Hongliang Zuo},
  title   = {Recent developments of matrix versions of the arithmetic--geometric mean inequality},
  journal = {Annals of Functional Analysis},
  volume  = {7},
  number  = {1},
  pages   = {102--117},
  year    = {2016}
}
\end{filecontents*}

\documentclass[11pt]{amsart}

%% ── wider margins ────────────────────────────────────────────────
\usepackage[margin=2.5cm]{geometry}

\usepackage[T1]{fontenc}
\usepackage[utf8]{inputenc}
\usepackage{amsmath,amssymb,amsthm,mathtools}
\usepackage{enumitem}
\usepackage{hyperref}
\usepackage[nameinlink,noabbrev]{cleveref}
\usepackage{xcolor}
\usepackage{microtype}
\usepackage{booktabs}
\usepackage{array}

\hypersetup{
  colorlinks=true,
  linkcolor=blue!60!black,
  citecolor=blue!60!black,
  urlcolor=blue!60!black
}

%% ── theorem environments ─────────────────────────────────────────
\newtheorem{theorem}{Theorem}[section]
\newtheorem{lemma}[theorem]{Lemma}
\newtheorem{proposition}[theorem]{Proposition}
\newtheorem{corollary}[theorem]{Corollary}
\theoremstyle{definition}
\newtheorem{definition}[theorem]{Definition}
\newtheorem{remark}[theorem]{Remark}
\newtheorem{example}[theorem]{Example}

%% ── macros ───────────────────────────────────────────────────────
\newcommand{\RR}{\mathrm{RR}}
\newcommand{\Sym}{\mathrm{Sym}}
\newcommand{\Wrd}{\mathrm{Wrd}}
\newcommand{\PSD}{\succeq 0}
\newcommand{\NSD}{\preceq 0}
\newcommand{\Loewner}{\preceq}
\newcommand{\Sn}{S_n}
\DeclarePairedDelimiter{\opnorm}{\lVert}{\rVert}
\newcommand{\tr}{\operatorname{tr}}
\newcommand{\diag}{\operatorname{diag}}
\newcommand{\R}{\mathbb{R}}
\newcommand{\C}{\mathbb{C}}
\newcommand{\N}{\mathbb{N}}
\newcommand{\Q}{\mathbb{Q}}
\newcommand{\Tr}{\operatorname{tr}}

%% ── title ────────────────────────────────────────────────────────
\title[Disproof of Bandeira's Problem 0.2(a)]{%
  Disproving Bandeira's Matrix AM--GM Open Problem~0.2(a):\\[4pt]
  An Exact Rational Certificate and Formal Verification}

\author{Anonymous}

\subjclass[2020]{15A45, 15A60, 47A63, 68V20, 90C22}
\keywords{Matrix AM--GM inequality, Recht--R\'{e} conjecture,
  semidefinite programming, Farkas certificate, formal verification, Lean~4}

\begin{document}

%% ═════════════════════════════════════════════════════════════════
\begin{abstract}
Bandeira's Open Problem~0.2(a) asks whether a matrix
arithmetic--geometric mean inequality holds in spectral norm for all
positive semidefinite matrices.  We give a complete mathematical
disproof by establishing that the problem is exactly the $m{=}n$
instance of the Recht--R\'e noncommutative AM--GM conjecture, and
invoking the Lai--Lim obstruction at $(m,n)=(5,5)$.

Our main contribution beyond the identification of the equivalence is
a fully reproducible verification chain.  Starting from Lai--Lim's
SDP data, we independently construct an exact rational Farkas-type
dual infeasibility certificate whose dual-vector entries have
denominator bound~$10^4$.  The $187{\times}187$ dual slack matrix
decomposes into one scalar coordinate and six independent
$31{\times}31$ blocks, each certified positive semidefinite by exact
rational LDL decomposition.  The entire finite
arithmetic---3\,906 constraint reconstructions, scalar contradiction,
six block LDL checks, and block-diagonal assembly---is
machine-checked in Lean~4 with Mathlib.  The only external
mathematical input is the equivalence chain established in
Lai--Lim~\cite[\S\,3]{LaiLim2020} (their Conjectures~1A--1C, proved
therein) from the Recht--R\'e inequality to the SDP standard form
via the noncommutative Positivstellensatz for bounded linear
matrix-inequality domains.

Supplementary material describes the Lean formalization in detail.
\end{abstract}

\maketitle
\tableofcontents

%% ═════════════════════════════════════════════════════════════════
\section{Introduction}\label{sec:intro}
%% ═════════════════════════════════════════════════════════════════

The classical arithmetic--geometric mean (AM--GM) inequality
\[
  \frac{1}{n}\sum_{i=1}^n a_i
  \;\ge\;
  \Bigl(\prod_{i=1}^n a_i\Bigr)^{1/n},
  \qquad a_i \ge 0,
\]
admits many matrix analogues, several of which have proven
surprisingly subtle and have driven substantial activity at the
interface of matrix analysis, free probability, and randomised
algorithms; see~\cite{Fujii2016} for a survey.  The particular
variant that motivates the present paper was posed by
Bandeira~\cite{Bandeira2015} as Open Problem~0.2 in his 2015 lecture
notes \emph{Ten Lectures and Forty-Two Open Problems in the
Mathematics of Data Science}.

For a tuple of positive semidefinite (PSD) matrices
$A_1,\dots,A_n \in \C^{d\times d}$ and the spectral (operator~$2$)
norm $\opnorm{\cdot}$, Bandeira asks whether the following two
inequalities always hold:
\begin{align}
  \label{eq:bandeira-a}
  \opnorm[\bigg]{\frac{1}{n!}\sum_{\sigma \in \Sn}
    \prod_{j=1}^n A_{\sigma(j)}}
    &\;\le\;
    \opnorm[\bigg]{\frac{A_1+\cdots+A_n}{n}}^{\!n},
    \tag{0.2a}\\[0.3em]
  \label{eq:bandeira-b}
  \frac{1}{n!}\sum_{\sigma \in \Sn}
    \opnorm[\big]{\prod_{j=1}^n A_{\sigma(j)}}
    &\;\le\;
    \frac{1}{n^n}\sum_{k_1,\dots,k_n=1}^{n}
    \opnorm[\big]{\prod_{j=1}^n A_{k_j}}.
    \tag{0.2b}
\end{align}
The cases $n=2$ are direct consequences of submultiplicativity and the
scalar AM--GM inequality.  The cases $n=3$ were established by
Israel, Krahmer, and Ward~\cite{IsraelKrahmerWard2016} via a
permutation variant of the Araki--Lieb--Thirring inequality, and
independently by Zhang.  Both inequalities have remained open for
general $n\ge 4$ since.

A separate but closely related line of work originates with Recht
and R\'e~\cite{RechtRe2012}, who introduced a noncommutative analogue
of AM--GM motivated by stochastic gradient descent (SGD) under
without-replacement sampling.  Their conjecture is the spectral-norm
inequality
\begin{equation}
  \label{eq:rrm}
  \frac{1}{n^{m}}
  \opnorm[\bigg]{\sum_{1\le j_1,\dots,j_m\le n}
    A_{j_1}\cdots A_{j_m}}
  \;\ge\;
  \frac{(n-m)!}{n!}
  \opnorm[\bigg]{\sum_{\substack{1\le j_1,\dots,j_m\le n\\
    j_1,\dots,j_m \;\text{distinct}}}
    A_{j_1}\cdots A_{j_m}},
  \tag{$\RR_{m,n}$}
\end{equation}
which compares the operator norm of the with-replacement matrix product
average to the without-replacement average.  The conjecture was
proven for $m=2$ and $m=3$ by Recht and R\'e, and strengthened by
Israel, Krahmer, and Ward~\cite{IsraelKrahmerWard2016} to general
unitarily invariant norms.  Lai and Lim~\cite{LaiLim2020}
subsequently proved the conjecture false by reducing it, via a
noncommutative Positivstellensatz~\cite{Helton2002}, to a semidefinite
program (SDP) and producing a dual infeasibility certificate at
$(m,n)=(5,5)$.

\subsection*{Contributions}
This paper makes the following contributions.
\begin{enumerate}[label=(\roman*),leftmargin=2em]
  \item We give a short, elementary proof
    (\cref{thm:equivalence}) that Bandeira's
    inequality~\eqref{eq:bandeira-a} for $n$ matrices is exactly the
    Recht--R\'e conjecture~\eqref{eq:rrm} at $m=n$ in spectral norm.
    The reduction relies on the identity
    $\opnorm{S^n}=\opnorm{S}^n$ for PSD~$S$.
  \item As a corollary (\cref{cor:main}), Bandeira's
    Problem~0.2(a) is \emph{false} for $n=5$: there exists a
    finite-dimensional PSD tuple
    $(A_1,\dots,A_5)$ for which~\eqref{eq:bandeira-a} fails.
  \item Starting from the Lai--Lim SDP data, we independently
    construct an \emph{exact rational} Farkas-type dual infeasibility
    certificate whose dual-vector entries have denominator bound~$10^4$.
    The $187{\times}187$ dual slack matrix decomposes into one scalar
    coordinate and six independent $31{\times}31$ blocks, each
    certified positive semidefinite by exact rational LDL decomposition
    (\cref{sec:exact}).
  \item We provide a machine-checked verification of the entire
    certificate in Lean~4 with Mathlib: no \texttt{sorry}, no
    \texttt{admit}, and no project-local axioms, covering the algebraic
    equivalence, the PSD spectral-norm power identity, the abstract
    Farkas contradiction, all 5\,766 standard-form block entries,
    six $31{\times}31$ exported PSD certificates, and the
    block-diagonal assembly into a single $(-M)\PSD$ theorem
    (\cref{sec:verification}).  All Lean details are deferred to
    the Supplementary Material.
  \item We identify the single formalization boundary: Lai--Lim's
    Section~3 equivalence chain (their Conjectures~1A--1C, proved
    therein) via the noncommutative Positivstellensatz is cited as a
    theorem from the literature (\cref{sec:boundary}).
  \item We disentangle Bandeira's parts~(a) and~(b), showing that
    the average-of-norms version~\eqref{eq:bandeira-b} is not
    equivalent to~\eqref{eq:rrm} and is not resolved by the present
    argument (\cref{sec:partb}).
\end{enumerate}

\smallskip
\noindent\textbf{Positioning.}\enspace
The counterexample is due to Lai and Lim~\cite{LaiLim2020}; our
contribution is the exact rational reconstruction of the certificate,
the machine-checked verification chain, and the precise separation of
the external reduction boundary.  This paper should be read as a
formal reconstruction and verification, not as a claim of a new
matrix counterexample.

\subsection*{Organisation}
\Cref{sec:setting} fixes notation and records the spectral-norm
statement of the Recht--R\'e conjecture.
\Cref{sec:equivalence} states and proves the equivalence with
Bandeira~0.2(a).
\Cref{sec:disproof} summarises the Lai--Lim reduction.
\Cref{sec:exact} presents the exact rational certificate and
its block structure.
\Cref{sec:verification} summarises the Lean verification
(details in Supplementary Material).
\Cref{sec:boundary} states the formalization boundary.
\Cref{sec:witnesses} discusses witness extraction.
\Cref{sec:partb} treats the companion inequality~0.2(b).

%% ═════════════════════════════════════════════════════════════════
\section{Setting and the spectral-norm Recht--R\'e conjecture}
\label{sec:setting}
%% ═════════════════════════════════════════════════════════════════

Throughout the paper, $\opnorm{\cdot}$ denotes the spectral
(operator~$2$) norm on matrices, $\PSD$ denotes positive
semidefiniteness, $\Loewner$ denotes the Loewner partial order on
Hermitian matrices, and $\Sn$ is the symmetric group on~$n$ elements.
All matrices are complex unless stated otherwise.

\begin{definition}[Symmetric and word averages]\label{def:averages}
Let $n,d\in\N$ and $A_1,\dots,A_n\in\C^{d\times d}$ with $A_i\PSD$.
For $1\le m\le n$, define
\begin{align}
  \Sym_m(A_1,\dots,A_n)
    &:= \frac{(n{-}m)!}{n!}
        \!\!\sum_{\substack{j_1,\dots,j_m\in[n]\\\text{distinct}}}
        \!\! A_{j_1}\cdots A_{j_m},
    \label{eq:sym}\\[4pt]
  \Wrd_m(A_1,\dots,A_n)
    &:= \frac{1}{n^m}
        \sum_{j_1,\dots,j_m=1}^{n} A_{j_1}\cdots A_{j_m}.
    \label{eq:wrd}
\end{align}
\end{definition}

\begin{remark}\label{rem:telescope}
For $m=n$, $\Sym_n$ averages over the $n!$ permutations and
$\Wrd_n$ averages over all $n^n$ length-$n$ words in
$\{1,\dots,n\}$.  The word sum telescopes:
$\sum_{j_1,\dots,j_n}A_{j_1}\cdots A_{j_n}=(A_1+\cdots+A_n)^{n}$,
hence $\Wrd_n(A_1,\dots,A_n)
= \bigl(\tfrac{A_1+\cdots+A_n}{n}\bigr)^{n}$.
\end{remark}

\begin{definition}[Recht--R\'e conjecture; spectral-norm form]
\label{def:rr}
Conjecture~$(\RR_{m,n})$ asserts that for every PSD tuple
$A_1,\dots,A_n$ in any finite dimension,
\[
  \opnorm{\Wrd_m(A_1,\dots,A_n)}
  \;\ge\;
  \opnorm{\Sym_m(A_1,\dots,A_n)}.
\]
\end{definition}

The original formulation of Recht and R\'e~\cite{RechtRe2012} uses the
Loewner order rather than the spectral norm; Lai and
Lim~\cite{LaiLim2020} record both versions and prove that the
conjecture in spectral norm is equivalent to a two-sided Loewner
statement (their Conjecture~1A) under the normalisation
$A_1+\cdots+A_n\Loewner nI$.

\begin{theorem}[{Recht--R\'e~\cite{RechtRe2012}; Israel--Krahmer--Ward~\cite{IsraelKrahmerWard2016}}; case $m\le 3$]
\label{thm:m23}
$(\RR_{m,n})$ holds for all $n$ when $m\le 3$.
\end{theorem}

\begin{theorem}[{Lai--Lim~\cite{LaiLim2020}}; main disproof]
\label{thm:lailim}
$(\RR_{m,n})$ is false at $m=n=5$.  Specifically, the smallest
constant $\lambda\ge 0$ for which the lower Loewner inequality
\[
  \sum_{\substack{j_1,\dots,j_5 \in [5] \\ \text{distinct}}}
  A_{j_1}\cdots A_{j_5} \;\succeq\; -\lambda\, I
\]
(under the normalisation $A_1+\cdots+A_5\Loewner 5I$) admits a
degree-$2$ noncommutative sum-of-squares certificate is
$\lambda_2 = 144.6488\ldots > 120 = 5!$.
\end{theorem}

\begin{remark}[Complex-to-real transfer]\label{rem:realification}
The equivalence in \cref{thm:equivalence} below is proved without
reference to the underlying field, so its conclusion holds over either
$\R$ or $\C$; when we appeal to Lai--Lim's disproof, we inherit their
setting, which is complex Hermitian.  Bandeira's original problem is
stated over real PSD matrices; the realification map
$\mathcal{R}(H+iK) :=
  \bigl(\begin{smallmatrix} H & -K \\ K & H \end{smallmatrix}\bigr)$
sends a complex Hermitian PSD matrix on $\C^d$ to a real symmetric
PSD matrix on $\R^{2d}$, preserves products
($\mathcal{R}(AB)=\mathcal{R}(A)\mathcal{R}(B)$), and preserves the
spectral norm (the singular values of $\mathcal{R}(A)$ are those
of~$A$, each with doubled multiplicity).  Consequently any finite-dimensional
complex PSD counterexample to~\eqref{eq:bandeira-a} yields a real PSD
counterexample of twice the dimension, and the disproof transfers to
Bandeira's original real-PSD formulation.
\end{remark}

%% ═════════════════════════════════════════════════════════════════
\section{The equivalence between Bandeira~0.2(a) and $(\RR_{n,n})$}
\label{sec:equivalence}
%% ═════════════════════════════════════════════════════════════════

The key observation is the following elementary identity.

\begin{lemma}[Norm of a PSD power]\label{lem:psdpower}
If $S\in\C^{d\times d}$ is positive semidefinite and $k\in\N$, then
$\opnorm{S^{k}} = \opnorm{S}^{k}$.
\end{lemma}

\begin{proof}
For any Hermitian matrix $M$,
$\opnorm{M}=\max_i \lvert\lambda_i(M)\rvert$.
For $S\PSD$, $\lambda_i(S)\ge 0$ and
$\lambda_i(S^{k})=\lambda_i(S)^{k}\ge 0$.  Hence
$\opnorm{S^{k}}
 =\max_i \lambda_i(S)^{k}
 =(\max_i \lambda_i(S))^{k}
 =\opnorm{S}^{k}$.
\end{proof}

\cref{lem:psdpower} fails without positivity: a nilpotent
matrix $N\ne 0$ has $\opnorm{N}>0$ but $\opnorm{N^{k}}=0$ for $k$
large enough.

\begin{lemma}[Word-sum identity]\label{lem:wordsum}
For elements $A_1,\dots,A_n$ in any (possibly noncommutative) semiring,
\[
  \sum_{j_1,\dots,j_n=1}^{n}
    A_{j_1}\cdots A_{j_n}
  = (A_1+\cdots+A_n)^n.
\]
\end{lemma}

\begin{proof}
Expand the power by the distributive law; each term corresponds to a
choice of indices $(j_1,\dots,j_n)\in[n]^n$.
\end{proof}

\begin{lemma}[Injection--permutation bijection]\label{lem:perm}
The set of injective functions $[n]\to[n]$ is exactly the set of
permutations $\Sn$, and its cardinality is $n!$.
\end{lemma}

\begin{theorem}[Main equivalence]\label{thm:equivalence}
For each $n\ge 1$, Bandeira's Open Problem~0.2(a) for $n$ matrices is
equivalent to $(\RR_{n,n})$.  More precisely, the inequalities
\eqref{eq:bandeira-a} and $(\RR_{n,n})$ in spectral norm are the same
inequality, for every PSD tuple $(A_1,\dots,A_n)$.
\end{theorem}

\begin{proof}
Fix $n\ge 1$ and $A_1,\dots,A_n\in\C^{d\times d}$ with $A_i\PSD$, and
set $S:=A_1+\cdots+A_n$.

\emph{Right-hand side of $(\RR_{n,n})$.}  The distinct-index
sum has exactly $n!$ terms, each corresponding to a permutation
$\sigma\in\Sn$ (\cref{lem:perm}); the coefficient is
$(n{-}n)!/n!=1/n!$.  Hence
\begin{equation}
  \frac{(n{-}n)!}{n!}
  \opnorm[\bigg]{\sum_{\text{distinct}}A_{j_1}\cdots A_{j_n}}
  = \frac{1}{n!}
    \opnorm[\bigg]{\sum_{\sigma\in\Sn}A_{\sigma(1)}\cdots A_{\sigma(n)}}.
  \label{eq:rhs}
\end{equation}

\emph{Left-hand side of $(\RR_{n,n})$.}  The word sum
telescopes (\cref{lem:wordsum}):
\[
  \sum_{j_1,\dots,j_n=1}^n A_{j_1}\cdots A_{j_n}
  = \prod_{r=1}^n \Bigl(\sum_{j_r=1}^n A_{j_r}\Bigr) = S^{n}.
\]
Since each $A_i\PSD$, $S\PSD$.  By \cref{lem:psdpower},
$\opnorm{S^{n}}=\opnorm{S}^{n}$, hence
\begin{equation}
  \frac{1}{n^{n}}
  \opnorm[\bigg]{\sum_{j_1,\dots,j_n}A_{j_1}\cdots A_{j_n}}
  = \frac{\opnorm{S}^n}{n^n}
  = \opnorm[\bigg]{\frac{A_1+\cdots+A_n}{n}}^{\!n}.
  \label{eq:lhs}
\end{equation}

\emph{Combining}~\eqref{eq:rhs} and~\eqref{eq:lhs}, the inequality
$(\RR_{n,n})$ at the tuple $(A_1,\dots,A_n)$ reads
\[
  \opnorm[\bigg]{\frac{A_1+\cdots+A_n}{n}}^{\!n}
  \;\ge\;
  \frac{1}{n!}
  \opnorm[\bigg]{\sum_{\sigma\in\Sn}
    A_{\sigma(1)}\cdots A_{\sigma(n)}},
\]
which after reversing the direction of inequality is
exactly~\eqref{eq:bandeira-a}.
\end{proof}

\begin{corollary}[Main result]\label{cor:main}
Bandeira's Open Problem~0.2(a) is \emph{false} for $n=5$.

Quantitatively, there exists a finite-dimensional PSD tuple
$(A_1,\dots,A_5)$ satisfying $A_1+\cdots+A_5 \Loewner 5I$ for which
the spectral-norm inequality~\eqref{eq:bandeira-a} fails.  The
Lai--Lim Loewner constant at $m=n=5$ is
$\lambda_2 = 144.6488\ldots > 120 = 5!$, which gives
$\lambda_2/5!\approx 1.2054$ as a strict upper bound on the smallest
constant $C$ for which a relaxed inequality could hold; we do not
claim a tuple attaining this gap.
\end{corollary}

\begin{proof}
By \cref{thm:equivalence}, \eqref{eq:bandeira-a} at $n=5$ is
equivalent to $(\RR_{5,5})$.  By \cref{thm:lailim}, $(\RR_{5,5})$ is
false.  Therefore Bandeira's inequality~\eqref{eq:bandeira-a} is false
for $n=5$.  The quantitative bound is read off from $\lambda_2/5!$.
\end{proof}

\begin{remark}[Status of the literature connection]\label{rem:literature}
The connection between Bandeira's Problem~0.2 and the Recht--R\'e
conjecture is recognised in the literature: the Lai--Lim paper is
widely cited in the random-reshuffling SGD literature, and Lai--Lim
themselves cite both Israel--Krahmer--Ward and Recht--R\'e.  We are
however not aware of a published statement of the \emph{exact}
equivalence in spectral norm at $m=n$ (\cref{thm:equivalence});
the literature describes the relationship variously as a
``special case'', ``related but distinct'', or as a partial overlap.
\Cref{lem:psdpower} renders it an identity, not merely a special case.
\end{remark}

%% ═════════════════════════════════════════════════════════════════
\section{The Lai--Lim reduction at $m=n=5$}
\label{sec:disproof}
%% ═════════════════════════════════════════════════════════════════

We summarise the structure of Lai--Lim's disproof, both for
self-containment and to fix the SDP that underlies the exact
certificate of \cref{sec:exact}.  The argument has four stages.

\paragraph{(i) Loewner reformulation.}
Lai and Lim~\cite[Conjecture~1A]{LaiLim2020} show that $(\RR_{m,n})$
in spectral norm is equivalent to the two-sided Loewner inequality
\begin{equation}
  -\frac{n!}{(n-m)!}\,I
  \;\Loewner\;
  \sum_{\substack{j_1,\dots,j_m\\\text{distinct}}}
    A_{j_1}\cdots A_{j_m}
  \;\Loewner\;
  \frac{n!}{(n-m)!}\,I,
  \label{eq:loewner}
\end{equation}
under the constraint $A_1+\cdots+A_n \Loewner nI$.  The upper bound
is the natural one; the lower bound is the one that fails at $m=n=5$.

\paragraph{(ii) Sum-of-squares certificate.}
Following the constrained noncommutative Positivstellensatz
framework~\cite{Helton2002,LaiLim2020}, the lower Loewner
inequality in~\eqref{eq:loewner}---i.e., that
$\sum_{\text{distinct}}A_{j_1}\cdots A_{j_m}\succeq -\lambda I$
under $\sum A_i\Loewner nI$---holds
(uniformly in dimension and matrix realisation) if and only if the
noncommutative polynomial
\[
  \lambda + \sum_{\text{distinct}} X_{j_1}\cdots X_{j_m}
\]
admits a constrained sum-of-squares representation
\[
  \lambda + \sum_{\text{distinct}} X_{j_1}\cdots X_{j_m}
  \;=\; \sigma_0 + \sum_{i=1}^{n+1} \sum_{l} h_{i,l}^{*}\, g_i\, h_{i,l},
\]
where $g_1,\dots,g_n := X_1,\dots,X_n$ are the PSD constraints,
$g_{n+1} := nI-\sum_j X_j$ encodes the normalisation, $\sigma_0$ is a
noncommutative sum of squares, and the total degree is at most
$2\lfloor m/2\rfloor+1$.

\paragraph{(iii) Semidefinite program.}
With the monomial basis $\beta$ of all noncommutative monomials of
degree $\le \lfloor m/2\rfloor$ in $X_1,\dots,X_n$ and the slack
variable $X_{n+1}:=n - \sum_i X_i$, matching coefficients
monomial-by-monomial yields the SDP of
\cite[Conjecture~1C]{LaiLim2020}.  At $m=n=5$ this SDP has six
$31\times 31$ moment matrices $Y_1,\dots,Y_6\PSD$.  The basis size is
\[
  \lvert\beta\rvert
  \;=\;\sum_{\ell=0}^{\lfloor m/2\rfloor} n^{\ell}
  \;=\;1+5+25\;=\;31,
\]
and the number of monomial-matching equality constraints is
\[
  \sum_{\ell=0}^{2\lfloor m/2\rfloor+1} n^{\ell}
  \;=\;1+5+25+125+625+3125\;=\;3906.
\]
The block-diagonal variable matrix has dimension
$1 + 6\cdot 31 = 187$; accounting for the symmetry of each
$31{\times}31$ block, the number of independent scalar
variables is $1 + 6\cdot 31\cdot 32/2 = 2977$, while the
total count of displayed block entries (used in the standard-form
reconstruction) is $6\cdot 31^{2} = 5766$.

\paragraph{(iv) Numerical verdict and Farkas certificate.}
Lai and Lim solve the SDP with SeDuMi~\cite{Sturm1999} and obtain
$\lambda_2\approx 144.6488$ at $m=n=5$.  To rule out numerical error,
they produce a Farkas dual certificate $y$ satisfying
$120\,y_0 + b^{\mathsf{T}}y > 0$, which by the Farkas lemma for
semidefinite programs~\cite{Lasserre2001} proves that the feasibility
problem at $\lambda=120$ is \emph{infeasible}, hence $\lambda_2>120$.

\begin{remark}[External boundary]\label{rem:external}
The reduction from \eqref{eq:loewner} to the SDP (stages
(i)--(iii) above) is the equivalence chain established in
\cite[\S\,3]{LaiLim2020} (their Conjectures~1A--1C, proved
therein), whose mathematical engine is the noncommutative
Positivstellensatz for bounded linear matrix-inequality
domains~\cite{Helton2002}.  This structural reduction---together
with the identification of the concrete standard-form coefficient
matrices $C_i$ and right-hand-side vector $b$ from Lai--Lim's
monomial-matching procedure---is the sole component of the
disproof that our Lean formalization does \emph{not} cover.
The identification of $C_i,b$ is reproduced and audited by
independent Python scripts (see \cref{sec:exact} and the
Supplementary Material).  We return to the formalization boundary
in \cref{sec:boundary}.
\end{remark}

%% ═════════════════════════════════════════════════════════════════
\section{Exact rational certificate}\label{sec:exact}
%% ═════════════════════════════════════════════════════════════════

Starting from the Lai--Lim SDP data, we independently construct an
exact rational Farkas-type dual infeasibility certificate.  This
section describes the mathematical structure; exact numerical values
are given in \cref{tab:cert}.

\subsection{Certificate structure}\label{sec:cert-structure}

A dual infeasibility certificate for the bound $\lambda=120$
consists of a vector $y=(y_0,\dots,y_{3905})\in\R^{3906}$
satisfying:
\begin{enumerate}[label=(\alph*),itemsep=2pt]
  \item \emph{Cone condition:} the dual slack matrix
    $M = y_0 C_0 + \sum_{i=1}^{3905} y_i C_i$
    satisfies $-M\PSD$.
  \item \emph{Scalar contradiction:}
    $120\,y_0 + b^{\mathsf{T}} y > 0$.
\end{enumerate}
Together these yield the contradiction
\[
  0
  \;<\; 120\,y_0 + b^{\mathsf{T}} y
  \;\le\; \langle M,Y\rangle
  \;\le\; 0,
\]
for any feasible $Y\PSD$, where the second inequality uses the
linear constraints $\langle C_i,Y\rangle = b_i$ and the last uses
$\tr(MY)\le 0$ for $M\NSD$ and $Y\PSD$.

\subsection{Rationalization and block decomposition}

Our construction proceeds as follows:
\begin{enumerate}[label=\arabic*.,itemsep=3pt]
  \item \emph{Numerics.}  We solve the Lai--Lim SDP independently
    using \texttt{cvxpy}+SCS, obtaining the optimal dual vector $z$
    with $\lambda^*\approx 144.649$.  We set $y = -z$ and
    $y_0 = -z_0$.

  \item \emph{Rationalization.}  Each $y_i$ is rounded to the nearest
    rational $p/q$ with $|q|\le 10^4$.  The resulting exact rational
    vector is denoted~$\hat{y}$.

  \item \emph{Block structure.}  The support graph of the constraint
    matrices $\{C_i\}$ has exactly seven connected components:
    one scalar coordinate ($m_{00}$) and six independent
    $31\times 31$ blocks.  The matrix $-M$ inherits this structure,
    with zero cross-block entries.

  \item \emph{Scalar verification.}  We verify in exact rational
    arithmetic:
    \[
      120\,\hat{y}_0 + b^{\mathsf{T}}\hat{y}
      = \frac{1\,417\,280}{1\,417\,281} > 0,
      \qquad
      m_{00} = -\frac{1\,060}{20\,438\,937} < 0.
    \]

  \item \emph{Block PSD certification.}  For each $31\times 31$ block
    $Q_b$ of $-M$, we compute the exact rational LDL decomposition
    $Q_b = L_b\diag(d_b)L_b^{\mathsf{T}}$ with all pivots
    $d_{b,k}>0$, certifying $Q_b\PSD$.
\end{enumerate}

\begin{table}[ht]
\centering
\caption{Key values of the exact rational certificate.  All have been
  verified in exact rational arithmetic and checked by Lean~4.}
\label{tab:cert}
\renewcommand{\arraystretch}{1.25}
\begin{tabular}{@{}ll@{}}
\toprule
\textbf{Quantity} & \textbf{Exact value} \\
\midrule
$y_0$ & $-146\,/\,2061$ \\[2pt]
$120\,y_0 + b^{\mathsf{T}} y$
  & $1\,417\,280\;/\;1\,417\,281 > 0$ \\[2pt]
$m_{00}$ (scalar block of $M$)
  & $-1\,060\;/\;20\,438\,937 < 0$ \\[2pt]
Block structure of $-M$
  & $1$ scalar $+$ $6$ blocks of $31{\times}31$ \\[2pt]
Each $31{\times}31$ block $Q_b$
  & LDL decomposition, all pivots ${>}\,0$ \\[2pt]
Denominator bound (entries of $\hat{y}$)
  & $10^4$ \\
\bottomrule
\end{tabular}
\end{table}

\subsection{Independent Python verification}

An independent Python verifier (using the \texttt{sympy} exact
rational library) reloads the certificate from disk and confirms:
(i)~the scalar positivity and negativity assertions;
(ii)~zero cross-block entries;
(iii)~exact LDL positive-definiteness for all six blocks.
The verifier runs in approximately 3~seconds and exits successfully.

%% ═════════════════════════════════════════════════════════════════
\section{Lean formalization overview}\label{sec:verification}
%% ═════════════════════════════════════════════════════════════════

The exact rational certificate obligations of \cref{sec:exact}, once
the standard-form matrices are fixed, and the algebraic equivalence of
\cref{sec:equivalence} have been verified in Lean~4~\cite{moura2021lean4}
with Mathlib~\cite{mathlib4}.  All theorem files compile with no
\texttt{sorry}, no \texttt{admit}, and no project-local axioms; only
standard Lean/Mathlib axioms are used.

\cref{tab:lean} summarises the formalization.  The detailed
file-by-file description, compilation commands, and the three-layer
split compilation strategy used for the $31\times 31$ blocks are
given in the Supplementary Material.

\begin{table}[ht]
\centering
\caption{Summary of the Lean~4 formalization.  ``Declarations''
  counts theorem and structure declarations.}
\label{tab:lean}
\renewcommand{\arraystretch}{1.25}
\begin{tabular}{@{}p{0.56\linewidth}rc@{}}
\toprule
\textbf{Component} & \textbf{Declarations} & \textbf{Status} \\
\midrule
Algebraic backbone
  (\cref{lem:wordsum,lem:perm,thm:equivalence})
                           & 14  & proved \\
PSD spectral-norm power (\cref{lem:psdpower})
                           &  9  & proved \\
Abstract Farkas chain      & 18  & proved \\
Standard-form reconstruction
                           & 191 & proved \\
LDL block checks ($6\times 64$)
                           & 384 & proved \\
Exported block PSD ($6\times 63+6$)
                           & 384 & proved \\
Block-diagonal assembly    &   3 & proved \\
\midrule
  \textbf{Total}
  & $\mathbf{>1{,}000}$
  & \textbf{clean} \\
\bottomrule
\end{tabular}
\end{table}

\noindent
The formalization has four layers:
\begin{enumerate}[label=(\roman*),leftmargin=2em,itemsep=3pt]
  \item \emph{Algebraic layer.}  Lean proofs of the word-sum identity,
    the injection--permutation bijection, the equivalence
    $\text{0.2(a)}\Leftrightarrow\RR_{n,n}$, and the integer bound
    $144>120$ (\texttt{ConcreteIdentities.lean}).  The PSD
    spectral-norm power identity is proved via a Hermitian
    conjugation identity, eigenvector preservation, and a
    characteristic polynomial bridge
    (\texttt{SpectralNormPSD.lean}).

  \item \emph{Abstract Farkas layer.}  Five files prove the scalar
    Farkas contradiction, the LDL-to-quadratic-form PSD bridge,
    the trace inner-product lemma $\tr(AB)\ge 0$ for $A,B\PSD$,
    an SDP infeasibility semantic bridge, and a Mathlib-native
    adapter connecting $(-M)\PSD$ to the scalar contradiction.

  \item \emph{Finite arithmetic layer.}  A generated file
    (\texttt{FarkasStandardFull.lean}, 191 theorems) verifies all
    5\,766 standard-form block entries by \texttt{norm\_num}.
    Six generated files (384 theorems total) verify the LDL
    row equalities and pivot positivity for all six $31\times 31$
    blocks.

  \item \emph{Exported PSD and assembly layer.}  For each of the six
    blocks, a three-layer split compilation pipeline
    (shared definitions $\to$ row arithmetic $\to$ row assembly
    $\to$ final assembly) proves $Q_b\PSD$ for the \emph{explicit
    exported} certificate matrices, not merely product-defined
    placeholders.  A final block-diagonal assembly file chains
    the scalar block and six $31\times 31$ PSD facts into a
    single theorem $(-M)_{\text{block-diag}}\PSD$ for the full
    $187\times 187$ dual slack matrix.
\end{enumerate}

%% ═════════════════════════════════════════════════════════════════
\section{Formalization boundary}\label{sec:boundary}
%% ═════════════════════════════════════════════════════════════════

All components described above---the algebraic equivalence, the
spectral-norm identity, the abstract Farkas chain, the finite
arithmetic, the six block PSD certificates, and the block-diagonal
assembly---are machine-checked in Lean with no \texttt{sorry},
no \texttt{admit}, and no project-local axioms (only Lean/Mathlib
standard axioms).

The \textbf{sole external input} is the equivalence chain
established in \cite[\S\,3]{LaiLim2020} (their
Conjectures~1A--1C, proved therein), which reduces the
Recht--R\'e inequality to the standard-form SDP via the
noncommutative Positivstellensatz for bounded linear
matrix-inequality domains~\cite{Helton2002}.
The identification of the concrete standard-form coefficient
matrices~$C_i$ and right-hand-side vector~$b$ with Lai--Lim's
monomial-matching SDP is reproduced by independent Python scripts
and verified by audit trails (see Supplementary Material), but is
not itself a Lean theorem.
We use the Lai--Lim reduction as a theorem from the literature;
once the SDP standard-form matrices are fixed, every numerical
component of the certificate---dual variables, scalar
contradiction, block decomposition, LDL pivots, and matrix
positive semidefiniteness---is verified by exact rational
arithmetic inside Lean.

%% ═════════════════════════════════════════════════════════════════
\section{Toward explicit witness matrices}\label{sec:witnesses}
%% ═════════════════════════════════════════════════════════════════

A natural strengthening of \cref{cor:main} would be to exhibit
explicit PSD matrices $A_1,\dots,A_5\in M_d(\C)$ that achieve the
violation.  Lai and Lim themselves note that ``it would be difficult
to find $A_1,\dots,A_5$ explicitly as one does not even know the
smallest dimensions required for these matrices to give a
counterexample''~\cite[\S\,3]{LaiLim2020}.

The standard route toward explicit witnesses is the
Gelfand--Naimark--Segal (GNS) construction.  In the moment-SDP
convention, an optimal or separating moment functional
$L\colon \R\langle X_1,\dots,X_5\rangle_{\le 5} \to \R$
that is positive on constrained sums of squares would be the
natural starting point.
The Hilbert space
$H := \R\langle X_1,\dots,X_5\rangle_{\le 2} / N_L$
(where $N_L$ is the kernel of the bilinear form
$\langle p,q\rangle = L(p^{*}q)$) is finite-dimensional, of dimension
at most~$31$.  Defining the candidate witnesses
$\pi(X_i)[p] := [X_i p]$ requires a \emph{flat-extension} condition
on~$L$: that the rank of the moment matrix stabilises at degree~$2$,
so that multiplication by $X_i$ descends to a well-defined
self-adjoint operator on~$H$.

The Farkas infeasibility certificate verified in the present
work proves that the feasibility problem at $\lambda=120$ has no
solution; it does not by itself provide the moment functional
required for witness extraction.  Obtaining such a functional
would require solving the explicit dual problem and verifying the
flat-extension condition, which we leave to future work.
We emphasise that the \emph{absence} of explicit witnesses does not
affect the truth of \cref{cor:main}: the disproof via
$\lambda_2 > 120$ is rigorous, and the existence of
finite-dimensional witnesses is guaranteed by the Positivstellensatz
\emph{a fortiori}.

%% ═════════════════════════════════════════════════════════════════
\section{The status of Bandeira's Problem~0.2(b)}\label{sec:partb}
%% ═════════════════════════════════════════════════════════════════

\Cref{thm:equivalence} concerns part~(a) only.  The companion
inequality~\eqref{eq:bandeira-b} places the norm inside the average:
\[
  \frac{1}{n!}\sum_{\sigma\in\Sn}
    \opnorm{A_{\sigma(1)}\cdots A_{\sigma(n)}}
  \;\le\;
  \frac{1}{n^n}\sum_{k_1,\dots,k_n=1}^{n}
    \opnorm{A_{k_1}\cdots A_{k_n}}.
\]
This is a \emph{scalar} inequality between two averages of
nonneg\-ative real numbers and is structurally different from the
matrix inequality~\eqref{eq:bandeira-a}.  It cannot be reduced to a
spectral-norm conjecture on a single PSD matrix-valued operator and
is therefore not covered by~$(\RR_{n,n})$.

\begin{lemma}[Commuting case]\label{lem:commuting}
If $A_1,\dots,A_n$ pairwise commute, then~\eqref{eq:bandeira-b} holds.
\end{lemma}

\begin{proof}
Commuting Hermitian operators admit a simultaneous spectral
decomposition $A_i = \sum_k \lambda_i^{(k)} e_k e_k^{*}$ with
$\lambda_i^{(k)}\ge 0$.  For a permutation $\sigma\in\Sn$,
\[
  \prod_{j=1}^{n} A_{\sigma(j)}
  =\sum_k\Bigl(\prod_{i=1}^{n}\lambda_i^{(k)}\Bigr)e_k e_k^{*},
\]
so the left side of~\eqref{eq:bandeira-b} equals
$L:=\max_k\prod_i\lambda_i^{(k)}$.  For a word
$w=(w_1,\dots,w_n)\in[n]^n$,
\[
  \opnorm[\big]{\prod_{r=1}^{n}A_{w_r}}
  =\max_k\prod_{r=1}^{n}\lambda_{w_r}^{(k)}.
\]
Choose $k_*$ attaining $L$.  Then
\begin{align*}
  \frac{1}{n^n}\sum_{w\in[n]^n}
    \max_k\prod_{r=1}^{n}\lambda_{w_r}^{(k)}
  &\ge
  \frac{1}{n^n}\sum_{w\in[n]^n}
    \prod_{r=1}^{n}\lambda_{w_r}^{(k_*)} \\
  &=\left(\frac{1}{n}\sum_{i=1}^{n}\lambda_i^{(k_*)}\right)^{\!n}
  \ge \prod_{i=1}^{n}\lambda_i^{(k_*)}=L,
\end{align*}
where the last inequality is the scalar AM--GM inequality.
\end{proof}

Any counterexample to~\eqref{eq:bandeira-b} at $n\ge 4$ must
therefore exploit a genuinely non-commuting configuration.
The present argument does not resolve part~(b); to our knowledge
it remains open.

\begin{remark}[Obstruction to a Positivstellensatz approach]
The spectral norm $\opnorm{A_{w_1}\cdots A_{w_n}}$ is not polynomial
in the matrix entries: $\opnorm{P}$ is the square root of the largest
eigenvalue of $P^{*}P$.  This is the obstruction to a direct
Lai--Lim-style Positivstellensatz attack on~\eqref{eq:bandeira-b}.
A natural substitute is the analogous statement for the Schatten
$p$-norm at finite even~$p$, where
$\opnorm{P}_{p}^{p}=\Tr\bigl((P^{*}P)^{p/2}\bigr)$ is a trace
polynomial; resolving the polynomial variant and taking the
$p\to\infty$ limit is a possible route.
\end{remark}

%% ═════════════════════════════════════════════════════════════════
\section{Conclusion}\label{sec:concl}
%% ═════════════════════════════════════════════════════════════════

We have shown that Bandeira's Open Problem~0.2(a) is identical to the
spectral-norm Recht--R\'e conjecture at $m=n$, via the elementary
identity $\opnorm{S^n}=\opnorm{S}^n$ for PSD~$S$, and is therefore
\emph{false} for $n=5$ by Lai and Lim's noncommutative
Positivstellensatz disproof.  The present argument does not resolve
the companion inequality~0.2(b); to our knowledge it remains open.

Our main contribution is the reproducible verification
infrastructure:
\begin{itemize}[itemsep=3pt]
  \item An exact rational Farkas-type dual infeasibility certificate
    (dual-vector denominator bound $10^4$), exploiting a
    seven-component block decomposition of the $187\times 187$ dual
    slack matrix.
  \item Machine-checked verification of all certificate arithmetic
    in Lean~4 with Mathlib, covering the algebraic equivalence,
    the spectral-norm power identity, the abstract Farkas chain,
    all standard-form reconstructions, six exported block PSD
    theorems, and the block-diagonal assembly
    (details in Supplementary Material).
  \item A single clearly identified external input: the equivalence
    chain established in Lai--Lim~\cite[\S\,3]{LaiLim2020}
    (Conjectures~1A--1C, proved therein) via the noncommutative
    Positivstellensatz.
\end{itemize}

The substantive remaining question is constructive: produce explicit
small-dimensional PSD matrices $A_1,\dots,A_5$ realising the
violation.  This requires solving the explicit noncommutative
moment-problem dual and verifying the flat-extension condition.

\subsection*{Code and data availability}

The source code, Lean~4 formalization, exact certificate data, and
verification scripts are organised in the accompanying repository
with the following structure:
\begin{center}
\renewcommand{\arraystretch}{1.15}
\begin{tabular}{@{}ll@{}}
\toprule
\textbf{Directory} & \textbf{Contents} \\
\midrule
\texttt{lean/}
  & Lean~4 proof files and \texttt{lakefile.lean} \\
\texttt{scripts/}
  & SDP reconstruction, certificate export, verification \\
\texttt{data/certificates/}
  & Exact rational certificate (\texttt{.npz}, \texttt{.json.gz}) \\
\texttt{paper/}
  & \LaTeX{} sources for this paper \\
\bottomrule
\end{tabular}
\end{center}

\bibliographystyle{alpha}
\bibliography{refs}

\end{document}
